\section{Standard scenario, held-out test, evasion and response}\label{sec:scenario}
\rubric{standard scenario \& evasion (20)}{We deliver the standard scenario, two frozen held-out sets with unseen methods, a robustness sweep, the trust-layer attack campaign, scenarios A--F and the response. Evidence: \cd{scripts/robustness_sweep.py}, \cd{scripts/heldout_v2.py}, \cd{scripts/run_attack_campaign.py}.}

\subsection{Standard scenario and held-out variants}
Authenticate clean patches, then detect a resampled or spliced image, a copy-move alteration and a generated or inpainted region, with held-out authentic and manipulated patches supplied by the jury. Our rehearsal uses manipulation variants \textbf{never seen in training} (the methods are listed in Table~\ref{tab:methods} of Section~\ref{sec:data}). Held-out v2 was created and SHA-256-frozen before any detector change made after v1; its fingerprint and the v1-model baseline on it are recorded, so the gain from our changes is honest.

\begin{table}[!htbp]
\centering\small
\begin{tabular}{|l|c|c|}
\hline
\textbf{Frozen held-out v2 (fused)} & \textbf{v1 model (baseline)} & \textbf{v2 model (shipped)}\\
\hline
detection rate & {{h1:/models/fused/detection_rate|pct1}} & \textbf{{{h:/models/fused/detection_rate|pct1}}}\\
\hline
false-alarm rate & {{h1:/models/fused/false_alarm_rate|pct1}} & \textbf{{{h:/models/fused/false_alarm_rate|pct1}}}\\
\hline
AUC-ROC & {{h1:/models/fused/auc_roc|f3}} & \textbf{{{h:/models/fused/auc_roc|f3}}}\\
\hline
copy-move & {{h1:/models/fused/detection_by_family/copy_move|pct1}} & \textbf{{{h:/models/fused/detection_by_family/copy_move|pct1}}}\\
\hline
splice & {{h1:/models/fused/detection_by_family/splice|pct1}} & \textbf{{{h:/models/fused/detection_by_family/splice|pct1}}}\\
\hline
resample & {{h1:/models/fused/detection_by_family/resample|pct1}} & \textbf{{{h:/models/fused/detection_by_family/resample|pct1}}}\\
\hline
inpaint & {{h1:/models/fused/detection_by_family/inpaint|pct1}} & \textbf{{{h:/models/fused/detection_by_family/inpaint|pct1}}}\\
\hline
\end{tabular}
\caption{Frozen held-out v2: the same set scored by the model before and after our changes. The false-alarm rate stays within the budget, so the detection gain is not bought with more false alarms.}
\end{table}

\textbf{Held-out v1, detection by family} (fused): resampling {{m:/models/fused (w=0.4)/test_heldout/detection_by_family/resample|pct1}}, inpainting/generation {{m:/models/fused (w=0.4)/test_heldout/detection_by_family/inpaint|pct1}}, copy-move {{m:/models/fused (w=0.4)/test_heldout/detection_by_family/copy_move|pct1}}, splicing {{m:/models/fused (w=0.4)/test_heldout/detection_by_family/splice|pct1}}.

\emph{Interpretation.} Detectors that read \emph{physical} traces generalise: interpolation leaves periodic correlations whatever the kernel, and every inpainting method leaves a region without sensor noise. Detectors that rely on \emph{shortcut} cues do not, which is where an adaptive attacker (A4) goes, and why provenance carries the guarantee.

\begin{figure}[!htbp]
\centering
\includegraphics[width=0.85\textwidth]{detection_by_family.png}
\caption{Detection rate by manipulation family and test condition (seen and held-out variants, with and without laundering).}
\end{figure}

\subsection{Our own red team: adversarial iterations}
Each iteration was evaluated on data never used for its development, and the measured failure drove the next design.
\begin{table}[!htbp]
\centering\small
\begin{tabularx}{\textwidth}{|L{0.6cm}|L{3.5cm}|L{3.6cm}|L{3.3cm}|Y|}
\hline
\textbf{\#} & \textbf{Problem found} & \textbf{Measured effect} & \textbf{Fix} & \textbf{Result}\\
\hline
1 & mild JPEG/noise augmentation only & \ext{14\%} false alarms on authentics after stronger re-encoding & widen augmentation & lower, but still leaking\\
\hline
2 & unseen re-encodings (WebP, blur) & \ext{45\%} of authentics looked manipulated & calibrate on clean validation & fewer, see \#3\\
\hline
3 & blur treated as benign & held-out resampling collapsed \ext{85\% $\to$ 33\%} (blur \emph{is} low-pass) & remove blur from the benign set & resampling restored\\
\hline
4 & re-encoded authentics still flagged & \ext{29--66\%} false alarms at JPEG q90--q45 & re-encoding gate and condition-aware threshold & see sweep below\\
\hline
5 & translation-only copy-move, hard-edge-only splice & held-out copy-move and splice weak & D4-invariant matching, colour-statistics features, feathered-edge augmentation & held-out v2 table above\\
\hline
\end{tabularx}
\caption{Adversarial iterations (development history; the iteration percentages are typed from the development log).}
\end{table}

\subsection{Robustness and the re-encoding gate}\label{sec:robust}
A detector trained on EuroSAT learns ``single-JPEG EuroSAT'', so a re-saved authentic patch can look manipulated. The gate (validation accuracy {{m:/gate_val_accuracy|pct1}}) moves the threshold between two calibrated points, keeping re-encoded authentics near the budget. The robustness sweep applies one operation to every one of {{r:/n_sources|int}} source patches:

\begin{table}[!htbp]
\centering\small
\begin{tabular}{|l|c|c|c|}
\hline
\textbf{Operation} & \textbf{Seen in training} & \textbf{False alarms} & \textbf{Detection}\\
\hline
none & yes & {{r:/rows/0/false_alarm_rate|pct1}} & {{r:/rows/0/detection_rate|pct1}}\\
\hline
JPEG q90 & yes & {{r:/rows/1/false_alarm_rate|pct1}} & {{r:/rows/1/detection_rate|pct1}}\\
\hline
JPEG q70 & yes & {{r:/rows/2/false_alarm_rate|pct1}} & {{r:/rows/2/detection_rate|pct1}}\\
\hline
JPEG q45 & yes & {{r:/rows/3/false_alarm_rate|pct1}} & {{r:/rows/3/detection_rate|pct1}}\\
\hline
JPEG q25 & no & {{r:/rows/4/false_alarm_rate|pct1}} & {{r:/rows/4/detection_rate|pct1}}\\
\hline
WebP q60 & yes & {{r:/rows/5/false_alarm_rate|pct1}} & {{r:/rows/5/detection_rate|pct1}}\\
\hline
blur 0.6 & no & {{r:/rows/6/false_alarm_rate|pct1}} & {{r:/rows/6/detection_rate|pct1}}\\
\hline
blur 1.2 & no & {{r:/rows/7/false_alarm_rate|pct1}} & {{r:/rows/7/detection_rate|pct1}}\\
\hline
noise 3 & yes & {{r:/rows/8/false_alarm_rate|pct1}} & {{r:/rows/8/detection_rate|pct1}}\\
\hline
noise 6 & no & {{r:/rows/9/false_alarm_rate|pct1}} & {{r:/rows/9/detection_rate|pct1}}\\
\hline
JPEG 2000 r10 & no & {{r:/rows/10/false_alarm_rate|pct1}} & {{r:/rows/10/detection_rate|pct1}}\\
\hline
JPEG 2000 r30 & no & {{r:/rows/11/false_alarm_rate|pct1}} & {{r:/rows/11/detection_rate|pct1}}\\
\hline
gamma 0.8 & no & {{r:/rows/12/false_alarm_rate|pct1}} & {{r:/rows/12/detection_rate|pct1}}\\
\hline
gamma 1.3 & no & {{r:/rows/13/false_alarm_rate|pct1}} & {{r:/rows/13/detection_rate|pct1}}\\
\hline
\end{tabular}
\caption{Robustness sweep, shipped detector. ``Seen in training'' means the operation was part of the laundering used in training.}
\end{table}

\begin{figure}[!htbp]
\centering
\includegraphics[width=0.95\textwidth]{robustness.png}
\caption{Effect of one operation applied to every patch (an asterisk marks an operation never seen in training); the dashed line is the false-alarm budget.}
\end{figure}

\begin{warnbox}
\textbf{Honest consequence.} Gaussian blur and JPEG 2000 on authentic patches still raise false alarms sharply, because they are physically the same low-pass signal as resampling. There is \textbf{no separate ``globally processed'' outcome yet}: a globally smoothed authentic patch is flagged like any manipulation (\cd{SIGNED_BUT_SUSPICIOUS} with a valid signature, \cd{SUSPICIOUS} without) and quarantined for an analyst. This costs a review, not a silently released fake. The integrity axis still proves, independently, that signed pixels are untouched since signing.
\end{warnbox}

\subsection{Trust-layer attack campaign}\label{sec:campaign}
{{c:/n_products|int}} genuine signed products, each attacked in transit:
\begin{table}[!htbp]
\centering\small
\begin{tabular}{|l|l|c|}
\hline
\textbf{Attack} & \textbf{Expected status} & \textbf{Caught}\\
\hline
none (genuine) & \cd{VERIFIED} & {{c:/attacks/0/caught_as_expected|pct0}}\\
\hline
inpaint after capture & \cd{PIXELS_ALTERED} & {{c:/attacks/1/caught_as_expected|pct0}}\\
\hline
copy-move after capture & \cd{PIXELS_ALTERED} & {{c:/attacks/2/caught_as_expected|pct0}}\\
\hline
resample after capture & \cd{PIXELS_ALTERED} & {{c:/attacks/3/caught_as_expected|pct0}}\\
\hline
LSB embedding (invisible) & \cd{PIXELS_ALTERED} & {{c:/attacks/4/caught_as_expected|pct0}}\\
\hline
date replaced, signature kept & \cd{BAD_SIGNATURE} & {{c:/attacks/5/caught_as_expected|pct0}}\\
\hline
date replaced and re-signed (attacker key) & \cd{UNKNOWN_KEY} & {{c:/attacks/6/caught_as_expected|pct0}}\\
\hline
hash downgraded to md5 & \cd{WEAK_ALGORITHM} & {{c:/attacks/7/caught_as_expected|pct0}}\\
\hline
forged processing step & \cd{CHAIN_BROKEN} & {{c:/attacks/8/caught_as_expected|pct0}}\\
\hline
manifest attached to another image & \cd{PIXELS_ALTERED} & {{c:/attacks/9/caught_as_expected|pct0}}\\
\hline
replay of a delivered product & \cd{REPLAY} & {{c:/attacks/10/caught_as_expected|pct0}}\\
\hline
edit + provenance stripped & \cd{NO_PROVENANCE} & {{c:/attacks/11/caught_as_expected|pct0}} (see note)\\
\hline
\end{tabular}
\caption{Trust-layer campaign: every attack against every product.}
\end{table}

\emph{Note on the last row.} ``{{c:/attacks/11/caught_as_expected|pct0}}'' means the verifier correctly reported that no provenance exists, \textbf{not} that the edit was stopped. Without a manifest only the content detector can notice: of {{c:/n_products|int}} edited patches whose manifest was removed, {{c:/two_axis_verdicts/strip_manifest/SUSPICIOUS|int}} were flagged \cd{SUSPICIOUS} and {{c:/two_axis_verdicts/strip_manifest/UNVERIFIED_NO_TRACES|int}} were released as \cd{UNVERIFIED_NO_TRACES}, i.e.\ not noticed.

\subsection{Attack scenarios A--F, run end to end}\label{sec:scenarios}
\cd{python -m simulation.run_attack --all --n {{at:/n_per_scenario|int}}} signs synthetic traffic from the bundled EuroSAT sample patches, attacks it, and feeds it to the real ground-station code (evidence for every trial in \cd{results/attack_runs/}). Every trial first delivers an untouched control product, which must be \cd{VERIFIED}, so a scenario can never pass because the ground station rejects everything. The first block is cryptographic: the expected status is a guarantee, and $n$ is the number of trials per scenario. The visual block has no signature to check, so the result is a measured detection rate, on manipulation methods the detector was not trained on, and it is small-sample.

\begin{table}[!htbp]
\centering\small
\begin{tabularx}{\textwidth}{|L{0.9cm}|Y|L{3.7cm}|L{3.3cm}|}
\hline
\textbf{ID} & \textbf{Scenario (what the attacker does)} & \textbf{Expected} & \textbf{Result}\\
\hline
A1 & image tampering: erase and fill a region, keep the manifest & \cd{PIXELS_ALTERED} & {{at:/scenarios/image_tampering/matches_expected|int}}/{{at:/n_per_scenario|int}}\\
\hline
A2 & invisible tampering: flip 512 least-significant bits & \cd{PIXELS_ALTERED} & {{at:/scenarios/invisible_tampering/matches_expected|int}}/{{at:/n_per_scenario|int}}\\
\hline
B1 & manifest tampering: change the capture date & \cd{BAD_SIGNATURE} & {{at:/scenarios/manifest_tampering/matches_expected|int}}/{{at:/n_per_scenario|int}}\\
\hline
B2 & identity forgery: claim to be the satellite, sign with an own key & \cd{UNKNOWN_KEY} & {{at:/scenarios/identity_forgery/matches_expected|int}}/{{at:/n_per_scenario|int}}\\
\hline
C1 & edit the image and rebuild a matching manifest, signed with an own key & \cd{BAD_SIGNATURE} & {{at:/scenarios/image_manifest_replacement/matches_expected|int}}/{{at:/n_per_scenario|int}}\\
\hline
C2 & algorithm substitution: claim md5 & \cd{WEAK_ALGORITHM} & {{at:/scenarios/algorithm_downgrade/matches_expected|int}}/{{at:/n_per_scenario|int}}\\
\hline
D1 & replay of a genuine, already delivered message & \cd{REPLAY} & {{at:/scenarios/replay/matches_expected|int}}/{{at:/n_per_scenario|int}}\\
\hline
E1 & copy-move, manifest removed & detector flags it & flagged {{at:/scenarios/visual_copy_move/content_flagged|int}}/{{at:/n_per_scenario|int}}\\
\hline
E2 & splicing, manifest removed & detector flags it & flagged {{at:/scenarios/visual_splice/content_flagged|int}}/{{at:/n_per_scenario|int}}\\
\hline
E3 & inpainting, manifest removed & detector flags it & flagged {{at:/scenarios/visual_inpaint/content_flagged|int}}/{{at:/n_per_scenario|int}}\\
\hline
E4 & resampling, manifest removed & detector flags it & flagged {{at:/scenarios/visual_resample/content_flagged|int}}/{{at:/n_per_scenario|int}}\\
\hline
F1 & fake made \textbf{before} signing (compromised payload) & integrity \cd{VERIFIED} & \cd{VERIFIED} {{at:/scenarios/compromised_source/matches_expected|int}}/{{at:/n_per_scenario|int}}; detector flagged {{at:/scenarios/compromised_source/content_flagged|int}}/{{at:/n_per_scenario|int}}\\
\hline
\end{tabularx}
\caption{Attack scenarios A--F run end to end.}
\end{table}

On the {{at:/n_per_scenario|int}} untouched originals used as controls, the detector raised {{at:/scenarios/visual_copy_move/control/false_alarms|int}} false alarms in E1 and {{at:/scenarios/visual_resample/control/false_alarms|int}} in E4. Scenario F1 is the honest limit of the whole design: a valid signature does not make a fake true.

\begin{figure}[!htbp]
\centering
\includegraphics[width=0.92\textwidth]{attack_lab.png}
\caption{The dashboard's Attack lab running scenario F1 (manipulated before signing): every fake carries a valid signature and the content detector flags only some of them. The panel shows what the attacker did, the expected and actual result, the untouched control, the evidence digests and the timeline of the satellite, the attacker and the ground station.}
\end{figure}

\subsection{Response}
Anything other than \cd{VERIFIED} or \cd{UNVERIFIED_NO_TRACES} is \textbf{quarantined}: kept, withheld downstream, released only by a logged analyst decision. The audit log is append-only and hash-chained; editing any past entry is detected (tested). The labs of Section~\ref{sec:labs} run the full satellite $\to$ attacker $\to$ ground chain and score it with \cd{simulation/score_demo.py}.
